\section{Frontier Workshop Notes}

\subsection{Workshop focus areas}

Lecturer 还建议定期 workshop，集中解决 offline RL 中的 open problems：

\begin{itemize}
    \item Lab session：Replay buffer introspection，手动审核 high-risk trajectories。
    \item Governance: simulate counterfactual audits + compliance docs in same meeting。
    \item Tooling：swap data versioning backends, evaluate time-to-recover metrics。
\end{itemize}

\begin{figure}[H]
\centering
\includegraphics[page=9,width=0.92\textwidth]{07_cs224r_offline_rl_2025.pdf}
\caption{Slides 第9页总结 workshop 议程与 checkpoint\protect\footnotemark}
\end{figure}
\footnotetext{Slides 第9页列出 workshop 的 checkpoint 及责任人。}

\subsection{本章小结}

Workshop 模式可以把 research insights 快速翻译成 operational improvements，避免 knowledge silo。

